\documentclass[12pt,a4paper]{article}

\usepackage[spanish,es-nodecimaldot]{babel}
\usepackage[utf8]{inputenc}
\usepackage[T1]{fontenc}
\usepackage{lmodern}
\usepackage{geometry}
\usepackage{microtype}
\usepackage{setspace}
\usepackage{parskip}
\usepackage{enumitem}
\usepackage{xcolor}
\usepackage{listings}
\usepackage{hyperref}
\usepackage{amsmath}
\usepackage{amssymb}
\usepackage{array}
\usepackage{booktabs}
\usepackage{longtable}
\usepackage{tcolorbox}

\geometry{margin=2.5cm}
\onehalfspacing
\setlist[itemize]{leftmargin=*,itemsep=0.2em}
\setlist[enumerate]{leftmargin=*,itemsep=0.2em}

\hypersetup{
  colorlinks=true,
  linkcolor=black,
  urlcolor=blue,
  pdftitle={Resumen teorico TP7 - Property-Based Testing y generacion aleatoria},
  pdfauthor={}
}

\newtcolorbox{definicionbox}[1][]{
  colback=blue!4!white,
  colframe=blue!50!black,
  fonttitle=\bfseries,
  title=Definicion,
  #1
}

\newtcolorbox{ideabox}[1][]{
  colback=green!4!white,
  colframe=green!40!black,
  fonttitle=\bfseries,
  title=Idea clave,
  #1
}

\newtcolorbox{ejemplobox}[1][]{
  colback=orange!4!white,
  colframe=orange!50!black,
  fonttitle=\bfseries,
  title=Ejemplo,
  #1
}

\lstdefinestyle{codestyle}{
  basicstyle=\ttfamily\small,
  frame=single,
  breaklines=true,
  showstringspaces=false
}

\title{\textbf{Resumen Teorico -- Practico 7}\\
Testing basado en propiedades y generacion aleatoria de tests\\
\textit{Testing de Software 2024}}
\author{}
\date{}

\begin{document}

\maketitle
\tableofcontents
\newpage

\section*{Introduccion}
\addcontentsline{toc}{section}{Introduccion}

Este resumen sintetiza el material teorico asociado al TP7:

\begin{itemize}
  \item \textit{Notas 12: Testing basado en propiedades} (V. Bengolea).
  \item \textit{Notas 13: Generacion automatica de tests --- generacion aleatoria guiada por
  feedback} (V. Bengolea).
\end{itemize}

\textbf{Idea general.} Hasta aqui los tests se escribian uno a uno: cada \texttt{@Test}
fijaba un input concreto y verificaba una salida concreta. Esa estrategia es valiosa pero
\emph{costosa, tediosa e incompleta}. Las dos tecnicas que se cubren en el TP7 amplian el
espectro:

\begin{itemize}
  \item \textbf{Property-Based Testing (PBT)}: en lugar de un input concreto, se enuncia una
  \emph{propiedad} que debe valer para todas las entradas dentro de una precondicion. El
  framework muestrea valores y reporta cualquier contraejemplo.
  \item \textbf{Generacion aleatoria automatica}: la herramienta no solo muestrea inputs, sino
  que construye \emph{secuencias de operaciones} aleatorias sobre las APIs del programa,
  alimentando \emph{feedback} de las ejecuciones anteriores para guiar la busqueda.
\end{itemize}

\section{Por que ir mas alla del unit testing}

\subsection{El espectro garantia vs.\ esfuerzo}

Las tecnicas de verificacion se pueden ordenar en un eje ``cuanta garantia dan'' vs.\ ``cuanto
esfuerzo cuestan'':

\begin{enumerate}
  \item \textbf{Metodos tradicionales} (unit tests escritos a mano): bajo esfuerzo por test,
  garantia muy local.
  \item \textbf{Testing automatico} (PBT, fuzzing, generacion aleatoria): mas garantia con
  esfuerzo medio. El programador solo escribe propiedades o configura generadores.
  \item \textbf{Metodos formales automaticos} (model checking, bounded verification): mayor
  garantia, costo de configuracion alto.
  \item \textbf{Verificacion deductiva} (Hoare, Coq, Isabelle): garantia maxima pero esfuerzo
  enorme.
\end{enumerate}

\begin{ideabox}
PBT y generacion automatica viven en el ``escalon util'' del espectro: dan mas garantia que
los unit tests sin pedir el esfuerzo de un metodo formal. Bien usados, pueden encontrar bugs
sutiles que ningun unit test individual cubriria.
\end{ideabox}

\subsection{Problemas de la creacion manual de tests}

Escribir tests a mano tiene tres problemas estructurales:

\begin{itemize}
  \item \textbf{Es costosa.} Cada caso requiere pensar input, oraculo y assertion.
  \item \textbf{Es tediosa.} Muchos casos son variaciones triviales.
  \item \textbf{Es incompleta.} El humano tiende a probar los casos que ya tiene en mente; los
  bordes y combinaciones inusuales suelen quedar fuera.
\end{itemize}

La automatizacion ataca los tres: testing automatico genera muchos casos rapidamente, y
explora regiones del input que un humano dificilmente cubriria.

\section{Property-Based Testing}

\subsection{Definicion}

\begin{definicionbox}
Un \textbf{test basado en propiedades} (Property-Based Test) captura un aspecto del
comportamiento del sistema bajo test para un \emph{numero potencialmente infinito de
entradas}. En lugar de un par fijo (input, output esperado), se enuncia una propiedad logica
que debe valer para \emph{todas} las entradas dentro de una precondicion.
\end{definicionbox}

\textbf{Contraste con unit testing.}

\begin{itemize}
  \item \textbf{Unit test}: comportamiento para una entrada particular.
  \item \textbf{Property test}: comportamiento para una familia entera de entradas.
\end{itemize}

\subsection{Estructura de una propiedad}

Una propiedad tipicamente sigue el esquema:

\begin{lstlisting}[style=codestyle]
public void myProperty(MyEntry e) {
    AssumePrecondition(e);    // 1. Precondicion: descartar entradas invalidas
    Arrange(e);               // 2. Preparar el escenario (opcional)
    Act(e);                   // 3. Ejecutar la operacion bajo test
    AssertExpectedProperty(e);// 4. Postcondicion: verificar la propiedad
}
\end{lstlisting}

Las cuatro partes:

\begin{itemize}
  \item \textbf{Precondicion (assume).} Describe las entradas validas. Si el generador produce
  algo que no la cumple, el caso se descarta (no se considera ``fail'').
  \item \textbf{Arrange.} Setup necesario antes de la operacion (e.g.\ poblar una lista).
  \item \textbf{Act.} La operacion bajo prueba.
  \item \textbf{Postcondicion (assert).} Lo que el codigo debe garantizar, expresado en
  funcion de los parametros que cumplen la precondicion.
\end{itemize}

\subsection{Ejemplo: \texttt{BoundedQueue.enqueue}}

\textbf{Propiedad.} Para toda cola no llena, insertar un elemento incrementa el tamano en 1.

\textbf{Unit test:}
\begin{lstlisting}[style=codestyle,language=Java]
@Test public void enqueueSizeTest_1() {
    BoundedQueue q = new BoundedQueue();
    q.enqueue(4); q.enqueue(3); q.enqueue(5);
    assertThat(q.size(), is(3));
}
\end{lstlisting}

Este test verifica la propiedad para un escenario puntual. Necesitariamos muchos tests
parecidos para distintos tamanos iniciales y elementos.

\textbf{Property test:}
\begin{lstlisting}[style=codestyle,language=Java]
public void enqueueSizeProperty(BoundedQueue queue, Integer i) {
    assumeTrue("cola no llena", !queue.isFull());
    int oldSize = queue.size();
    queue.enqueue(i);
    assertThat(queue.size(), is(oldSize + 1));
}
\end{lstlisting}

Una sola propiedad reemplaza a una familia entera de unit tests, y el framework muestrea
muchos valores de \texttt{queue} y de \texttt{i}.

\subsection{Mas ejemplos clasicos de propiedades}

\begin{itemize}
  \item \textbf{Sets:} ``si se agrega un elemento que ya pertenece, el set no cambia''.
  $$\forall s, x.\ x \in s \Rightarrow s \cup \{x\} = s.$$
  \item \textbf{Strings:} ``$a$ es substring de \texttt{concat}$(a, b)$ para todo $a, b$''.
  \item \textbf{Math.max:} para todo $a, b$ enteros,
  \begin{itemize}
    \item \texttt{Math.max(a, b) == Math.max(b, a)} (conmutatividad),
    \item \texttt{Math.max(a, b) >= a} y \texttt{>= b},
    \item \texttt{p + Math.max(a, b) == Math.max(p+a, p+b)} (compatibilidad con suma).
  \end{itemize}
\end{itemize}

\subsection{Tipos de propiedades utiles}

\begin{itemize}
  \item \textbf{Invariantes.} Propiedades que se preservan al aplicar operaciones (e.g.\
  \texttt{repOK()} antes y despues).
  \item \textbf{Contratos del lenguaje.} Como el contrato \texttt{equals/hashCode} de
  \texttt{Object}: si \texttt{a.equals(b)} entonces \texttt{a.hashCode() == b.hashCode()}.
  \item \textbf{Propiedades algebraicas.} Conmutatividad, asociatividad, distributividad,
  idempotencia.
  \item \textbf{Round-trip.} \texttt{deserialize(serialize(x)) == x}, \texttt{parse(print(t))
  == t}, \texttt{decompress(compress(d)) == d}.
  \item \textbf{Oraculo de referencia.} Si existe una implementacion lenta pero confiable, la
  propiedad puede ser ``mi implementacion rapida da el mismo resultado que la lenta''.
  \item \textbf{Metamorficas.} Relacion entre dos llamadas del mismo metodo:
  \texttt{sort(reverse(l)) == sort(l)}.
\end{itemize}

\subsection{Como hace el framework para muestrear}

Un framework PBT (e.g.\ \texttt{jqwik} en Java, \texttt{QuickCheck} en Haskell,
\texttt{Hypothesis} en Python) ejecuta cada propiedad cientos de veces. En cada iteracion:

\begin{enumerate}
  \item Llama a los generadores (\texttt{@ForAll}/\texttt{@Provide}) para producir valores
  para los parametros.
  \item Aplica la precondicion (\texttt{assumeTrue}); si falla, descarta el caso y prueba
  otro.
  \item Ejecuta la propiedad. Si lanza una excepcion o falla una assertion, registra el
  contraejemplo.
\end{enumerate}

\subsection{Shrinking: minimizar contraejemplos}

\begin{definicionbox}
\textbf{Shrinking.} Cuando una propiedad falla, el framework intenta \emph{simplificar} el
input que provoco la falla para obtener el contraejemplo \emph{minimo}. Por ejemplo, si una
lista de 100 elementos rompe la propiedad, el shrinking busca la sub-lista mas chica que aun
la rompe.
\end{definicionbox}

\begin{ideabox}
Un contraejemplo minimo es muchisimo mas facil de depurar que uno gigante con datos al azar.
Por eso shrinking es lo que distingue a un framework PBT serio de un simple ``ejecuta esto
mil veces''.
\end{ideabox}

\section{Generadores}

\subsection{Por que importan tanto}

Las propiedades son tan buenas como los generadores. Un generador que solo produce listas
vacias hace que la propiedad ``ordenar preserva la longitud'' sea trivialmente verdadera. Hay
que pensar:

\begin{itemize}
  \item Que entradas son \emph{validas}.
  \item Que entradas son \emph{representativas}.
  \item Que entradas son \emph{borde} (valores extremos, vacios, repetidos).
\end{itemize}

\subsection{Primitivos: enteros, floats, strings}

Casi todos los frameworks proveen primitivos con rangos configurables:

\begin{lstlisting}[style=codestyle,language=Java]
Arbitraries.integers().between(0, 100);
Arbitraries.floats().between(-1e3f, 1e3f);
Arbitraries.strings().withCharRange('a', 'z').ofMaxLength(20);
\end{lstlisting}

\subsection{Combinadores}

A partir de primitivos se arman generadores compuestos. Los dos mas importantes son
\texttt{map} y \texttt{flatMap}:

\begin{itemize}
  \item \textbf{\texttt{map}}: aplica una funcion al valor generado.
  \begin{lstlisting}[style=codestyle,language=Java]
  Arbitraries.integers().between(0, 100).map(n -> "user" + n);
  \end{lstlisting}
  \item \textbf{\texttt{flatMap}}: permite que un parametro \emph{dependa} de otro.
  \begin{lstlisting}[style=codestyle,language=Java]
  // generador de listas y un indice valido sobre cada lista
  Arbitraries.integers().between(1, 100).flatMap(size ->
    Arbitraries.integers().list().ofSize(size).flatMap(list ->
      Arbitraries.integers().between(0, size - 1).map(idx ->
        new Scenario(list, idx))));
  \end{lstlisting}
\end{itemize}

\subsection{Generadores dependientes y validez}

Muchos dominios requieren \emph{validez compuesta}: un dia depende del mes y del ano, un
indice depende del tamano de la lista, una transicion de FSM depende del estado actual.
\texttt{flatMap} es la herramienta para construir esos generadores sin tener que filtrar
muchas entradas invalidas.

\begin{ejemplobox}
Generar fechas validas a partir de 1900:
\begin{lstlisting}[style=codestyle,language=Java]
Arbitraries.integers().between(1900, 2400).flatMap(year ->
  Arbitraries.integers().between(1, 12).flatMap(month ->
    Arbitraries.integers().between(1, daysInMonth(month, year))
      .map(day -> new Date(day, month, year))));
\end{lstlisting}
Sin \texttt{flatMap} se generarian fechas como \texttt{31/04/2024}, que el constructor
rechazaria.
\end{ejemplobox}

\subsection{\texttt{filter} vs.\ \texttt{flatMap}}

\texttt{filter} descarta valores que no cumplen una condicion. Es \emph{lo ultimo} a usar:
si la condicion es estricta (por ejemplo, ``primos''), el framework descartara miles de
muestras hasta encontrar una valida.

\begin{ideabox}
Regla practica: si el filtro descarta mas del $\sim$30\% de las muestras, conviene rediseniar
el generador con \texttt{flatMap} o construir directamente valores validos.
\end{ideabox}

\section{Generacion automatica de tests}

\subsection{Idea general}

La generacion aleatoria automatica de tests va un paso mas alla del PBT: en vez de pedirle al
programador que escriba propiedades y generadores, la herramienta:

\begin{enumerate}
  \item Recibe el codigo del programa (sus clases, metodos, constructores).
  \item Genera \emph{secuencias de llamadas} aleatorias a esas APIs.
  \item Ejecuta cada secuencia y observa que pasa (excepciones, valores devueltos, contratos
  violados).
  \item Usa esa informacion como \emph{feedback} para decidir como continuar.
\end{enumerate}

\subsection{Generacion aleatoria guiada por feedback}

\begin{definicionbox}
\textbf{Generacion guiada por feedback} (e.g.\ \emph{Randoop}). Se genera la secuencia de
llamadas paso a paso, ejecutando cada paso inmediatamente. Si la ejecucion lanza una
excepcion no contractual (e.g.\ un \texttt{NullPointerException} no esperado), la secuencia
se reporta como \emph{test que muestra un bug}. Si la ejecucion va bien, el resultado se
guarda en un \emph{pool} de valores que puede usarse para construir secuencias futuras.
\end{definicionbox}

\textbf{Ciclo basico:}

\begin{enumerate}
  \item Comenzar con un pool de valores trivial (enteros, strings).
  \item Elegir un metodo al azar de la API bajo prueba.
  \item Elegir, del pool, valores que sirvan como argumentos compatibles en tipos.
  \item Ejecutar la llamada:
  \begin{itemize}
    \item Si la ejecucion falla con un crash inesperado, reportar la secuencia como test
    sospechoso.
    \item Si retorna un valor, agregarlo al pool.
    \item Si lanza una excepcion documentada, marcar como ``no extendible'' y descartar.
  \end{itemize}
  \item Repetir hasta agotar el presupuesto de tiempo.
\end{enumerate}

\subsection{Oraculo}

Un detalle clave: \emph{quien decide} si una ejecucion es un bug? Hay varios oraculos posibles:

\begin{itemize}
  \item \textbf{Crashes inesperados.} \texttt{NullPointerException},
  \texttt{ArrayIndexOutOfBoundsException}, \texttt{AssertionError} en lugares donde la
  especificacion no los permite.
  \item \textbf{Contratos del lenguaje.}
  \texttt{equals}/\texttt{hashCode}/\texttt{toString}/\texttt{compareTo} con sus reglas:
  reflexividad, simetria, transitividad, etc.
  \item \textbf{Aserciones internas} del codigo (\texttt{assert} de Java).
  \item \textbf{Diferencial.} Comparar contra una version de referencia.
\end{itemize}

\subsection{Por que feedback}

La aleatoriedad pura genera muchas secuencias \emph{redundantes} (mismas operaciones, mismos
valores). La guia por feedback hace varias cosas:

\begin{itemize}
  \item Evita repetir secuencias equivalentes (memoization de pools).
  \item Acumula objetos en estados \emph{interesantes} que luego pueden combinarse.
  \item Detecta secuencias que terminan rapido (excepciones esperadas) y las descarta como
  ``muertas''.
\end{itemize}

\subsection{Limitaciones}

\begin{itemize}
  \item \textbf{Oraculo limitado.} Sin propiedades explicitas, la herramienta solo detecta
  bugs ``visibles'' (crashes, contratos de lenguaje). Las violaciones de logica de dominio
  pasan inadvertidas.
  \item \textbf{Tests poco legibles.} Los tests generados son secuencias mecanicas; cuando
  fallan, son dificiles de entender.
  \item \textbf{Cobertura sesgada.} La aleatoriedad uniforme no es exploratoria: ciertas
  partes del codigo quedan profundamente cubiertas y otras no.
\end{itemize}

\section{Relacion entre PBT, fuzzing y generacion automatica}

\begin{center}
\begin{tabular}{lll}
\toprule
Tecnica & Quien escribe que & Que entrega \\
\midrule
Unit testing & Programador escribe input y oraculo & Tests concretos \\
PBT & Programador escribe propiedad y generador & Familia de tests \\
Fuzzing & Programador da semillas u oraculo minimo & Inputs (rara vez tests legibles) \\
Generacion automatica & Programador da API (o nada) & Secuencias de llamadas \\
\bottomrule
\end{tabular}
\end{center}

\begin{ideabox}
A medida que se baja en la tabla, el programador escribe menos pero \emph{el oraculo se
debilita}. PBT es un buen punto medio: la propiedad da un oraculo fuerte (no solo ``no
crashea''), y el generador da exploracion amplia.
\end{ideabox}

\section{Patrones para disenar propiedades}

\subsection{El patron ``oracle por construccion''}

Cuando es dificil enunciar el resultado correcto, a veces se puede construir el escenario
\emph{con} el resultado ya conocido:

\begin{ejemplobox}
En vez de generar un grafo y verificar ``el camino mas corto entre $u$ y $v$'', se construye
un grafo donde ya se sabe el camino mas corto (por ejemplo, un grafo con un solo camino
posible) y se verifica que el algoritmo lo encuentra.
\end{ejemplobox}

\subsection{El patron ``operaciones inversas''}

Si una operacion tiene una inversa (deserializar/serializar, encriptar/desencriptar,
parse/print, push/pop), la propiedad de round-trip es muy efectiva: aplica ambas y verifica
que se vuelve al estado inicial.

\subsection{El patron ``invariante de representacion''}

Pedir que \texttt{repOK()} se preserve por toda operacion publica. Captura una clase entera
de bugs (corrupcion de estado interno) con una propiedad sola.

\subsection{El patron ``ley algebraica''}

Si el dominio tiene leyes (asociatividad, conmutatividad, identidad), enunciar cada ley como
una propiedad.

\section{Comparacion con el resto del curso}

\begin{itemize}
  \item \textbf{Particion del input} (TP2) y PBT son complementarias: PBT muestrea \emph{en}
  cada bloque, pero la particion ayuda a definir los rangos del generador.
  \item \textbf{Cobertura estructural} (TP3-TP4) responde a una pregunta distinta:
  ``\textquestiondown{}toque todo el codigo?''. PBT responde ``\textquestiondown{}mi codigo
  cumple la propiedad?''. Idealmente ambas se combinan: usar cobertura para asegurarse de que
  los generadores ejercitan todas las ramas.
  \item \textbf{Cobertura logica} (TP5) y PBT se relacionan: las propiedades suelen escribirse
  como predicados, y los generadores pueden disenarse para ejercitar CACC de esos predicados.
  \item \textbf{Mutation testing} (TP6) puede medir la calidad de una suite de propiedades:
  un buen conjunto de propiedades deberia matar la mayoria de los mutantes igual que una
  buena suite de unit tests.
  \item \textbf{Fuzzing} (TP6) es generacion aleatoria \emph{sin} propiedades sofisticadas:
  el oraculo es ``no crash''. PBT y generacion guiada por feedback agregan oraculos mas
  precisos.
\end{itemize}

\section{Glosario}

\begin{itemize}
  \item \textbf{Propiedad}: enunciado que debe valer para todas las entradas que cumplan una
  precondicion.
  \item \textbf{Generador} (\emph{arbitrary}, \emph{provider}): mecanismo que produce valores
  aleatorios para un parametro.
  \item \textbf{Shrinking}: proceso de minimizar un contraejemplo cuando una propiedad falla.
  \item \textbf{Precondicion}: condicion sobre las entradas para que la propiedad sea
  aplicable.
  \item \textbf{Postcondicion}: lo que el codigo debe garantizar.
  \item \textbf{Oraculo}: mecanismo que decide si la ejecucion es correcta.
  \item \textbf{Pool de valores}: conjunto de objetos disponibles para construir nuevas
  secuencias en generacion automatica.
\end{itemize}

\section*{Conclusion}
\addcontentsline{toc}{section}{Conclusion}

PBT cambia la unidad de testing: en vez de pensar en \emph{un} input, se piensa en una
\emph{familia} entera de inputs. Esa abstraccion fuerza al programador a hacer explicitas
las propiedades que su codigo realmente debe satisfacer --- y muchas veces, al hacerlo, se
descubre que las propiedades no estaban tan claras. La generacion automatica de tests
empuja aun mas la automatizacion: con una API y un oraculo minimo, la herramienta busca
secuencias que crasheen. Ninguna de las dos reemplaza al unit testing, pero ambas atacan
problemas que el unit testing solo no resuelve: cobertura sesgada, casos borde
no-pensados y bugs combinatorios.

\end{document}
